\documentclass[12pt]{article}
\usepackage[ukrainian]{babel}
\usepackage{fontspec}
% --- НАЛАШТУВАННЯ ШРИФТІВ ---
\setmainfont{Schoolbook-Regular.otf}[
    Path = ../../,
    BoldFont = Schoolbook-Bold.otf,
    ItalicFont = Schoolbook-Italic.otf,
    BoldItalicFont = Schoolbook-BoldItalic.otf
]
\usepackage[italic]{mathastext}
\usepackage[a4paper,margin=1.8cm,bottom=2cm,top=2cm]{geometry}
\usepackage{amsmath,amssymb}
\usepackage{tikz}
\usepackage{pgfplots}
\pgfplotsset{compat=1.16}
\usetikzlibrary{calc,patterns,angles,quotes}
\usepackage{xcolor,array,fancyhdr,enumitem,multicol}
\usepackage{colortbl}
\usepackage{tcolorbox}
\tcbuselibrary{skins, breakable}
\usepackage[hidelinks]{hyperref}
\usepackage[firstpage=false, text=@pvtr2525, scale=0.6, color=gray!12]{draftwatermark}
% --- ФІРМОВІ КОЛЬОРИ ---
\definecolor{mainGreen}{RGB}{34, 120, 64}
\definecolor{yearOrange}{RGB}{220, 100, 30}
\definecolor{yearcolor}{RGB}{220, 100, 30}
\definecolor{titleBg}{RGB}{225, 240, 220}
\definecolor{instrBg}{RGB}{255, 240, 220}
\definecolor{instrBorder}{RGB}{220, 150, 80}
\pagestyle{fancy}
\fancyhf{}
\renewcommand{\headrulewidth}{0.4pt}
\renewcommand{\footrulewidth}{0.4pt}
\fancyhead[C]{\small\color{gray!80} Збірник НМТ \textendash{} сесії 23.05--5.06.2026}
\fancyhead[R]{\small\color{mainGreen}\textbf{\href{https://t.me/pvtr2525}{@pvtr2525}}}
\fancyfoot[L]{\small\color{gray!80}\href{https://t.me/pvtr2525}{ТГ: @pvtr2525}}
\fancyfoot[C]{\small\color{gray!80}\thepage}
\fancyfoot[R]{\small\color{gray!80}НМТ 2026}
% --- ЗАГОЛОВОК РОЗДІЛУ ---
\newcommand{\sectionTitle}[1]{%
\par\vspace{0.4cm}
\begin{tcolorbox}[colback=titleBg, colframe=titleBg, boxrule=0pt, arc=8pt,
    halign=center, fontupper=\Large\bfseries\color{mainGreen}]
#1
\end{tcolorbox}
\par\vspace{0.3cm}
}
% --- МАКРОСИ З ГІПЕРПОСИЛАННЯМИ ---
\newcommand{\answerTable}[5]{%
\par\vspace{0.15cm}
\noindent
\begin{tabular}{|*{5}{>{\centering\arraybackslash}m{3cm}|}}
\hline
\rule[-0.15cm]{0pt}{0.6cm}\textbf{А} & \rule[-0.15cm]{0pt}{0.6cm}\textbf{Б} & \rule[-0.15cm]{0pt}{0.6cm}\textbf{В} & \rule[-0.15cm]{0pt}{0.6cm}\textbf{Г} & \rule[-0.15cm]{0pt}{0.6cm}\textbf{Д} \\
\hline
\rule[-0.2cm]{0pt}{0.75cm}#1 & \rule[-0.2cm]{0pt}{0.75cm}#2 & \rule[-0.2cm]{0pt}{0.75cm}#3 & \rule[-0.2cm]{0pt}{0.75cm}#4 & \rule[-0.2cm]{0pt}{0.75cm}#5 \\
\hline
\end{tabular}
\par\vspace{0.25cm}
}
\newcommand{\answerTableTall}[5]{%
\par\vspace{0.15cm}
\noindent
\begin{tabular}{|*{5}{>{\centering\arraybackslash}m{3cm}|}}
\hline
\rule[-0.2cm]{0pt}{0.7cm}\textbf{А} & \rule[-0.2cm]{0pt}{0.7cm}\textbf{Б} & \rule[-0.2cm]{0pt}{0.7cm}\textbf{В} & \rule[-0.2cm]{0pt}{0.7cm}\textbf{Г} & \rule[-0.2cm]{0pt}{0.7cm}\textbf{Д} \\
\hline
\rule[-0.45cm]{0pt}{1.2cm}#1 & \rule[-0.45cm]{0pt}{1.2cm}#2 & \rule[-0.45cm]{0pt}{1.2cm}#3 & \rule[-0.45cm]{0pt}{1.2cm}#4 & \rule[-0.45cm]{0pt}{1.2cm}#5 \\
\hline
\end{tabular}
\par\vspace{0.25cm}
}
\newcommand{\taskBlock}[1]{%
\par\vspace{0.45cm}
\noindent{\color{mainGreen}\rule{\linewidth}{0.8pt}}\par\vspace{0.12cm}
\noindent{\small\bfseries\color{mainGreen}#1}\par\vspace{0.25cm}
}
% --- ГІПЕРПОСИЛАННЯ В НОМЕРАХ ЗАВДАНЬ ---
\newcounter{zad}
\newcommand{\zadtask}[1]{\stepcounter{zad}\noindent\hypertarget{zad\thezad}{}\hyperlink{ans\thezad}{\textbf{\thezad.}}\ \ #1\par\vspace{0.2cm}}
\newcommand{\zadnum}{\stepcounter{zad}\hypertarget{zad\thezad}{}\hyperlink{ans\thezad}{\textbf{\thezad.}}\ \ }
\newcounter{ans}
\newcommand{\ansitem}[1]{\stepcounter{ans}\noindent\hypertarget{ans\theans}{}\hyperlink{zad\theans}{\textbf{\theans.}}\ #1\par\vspace{0.07cm}}
\newcommand{\nmtyear}[1]{\hfill{\small\color{yearOrange}(НМТ #1)}}
\newcommand{\ansTheme}[1]{\par\vspace{0.25cm}\noindent{\bfseries\color{mainGreen}#1}\par\vspace{0.12cm}}
\newcommand{\ansType}[1]{\par\vspace{0.3cm}\noindent{\large\bfseries\color{yearOrange}#1}\par\vspace{0.15cm}}
% --- ПРИХОВАНІ РОЗВ'ЯЗКИ З КЛІКАБЕЛЬНИМИ ПОСИЛАННЯМИ ---
\newcounter{solctr}
\makeatletter
\newcommand{\solution}[1]{%
\stepcounter{solctr}%
\hypertarget{task:\thesolctr}{}%
\par\nobreak\vspace{0.05cm}\hfill%
\hyperlink{sol:\thesolctr}{\small\bfseries\color{mainGreen}\fbox{\strut\,розв'язок $\rightarrow$\,}}%
\par\vspace{0.15cm}%
\expandafter\xdef\csname ztasknum@\thesolctr\endcsname{\thezad}%
\expandafter\gdef\csname zsol@\thesolctr\endcsname{#1}%
}
\newcommand{\printAllSolutions}{%
\chapterTitle{РОЗВ'ЯЗКИ ОРИГІНАЛЬНИХ ЗАДАЧ}%
\noindent\small Натисни на номер задачі ліворуч від рамки, щоб повернутися до умови.\par\vspace{0.3cm}%
\newcounter{psctr}%
\@whilenum\value{psctr}<\value{solctr}\do{%
\stepcounter{psctr}%
\par\noindent\hypertarget{sol:\thepsctr}{}%
\begin{tcolorbox}[colback=titleBg!45, colframe=mainGreen!55, boxrule=0.5pt, arc=3pt,
    left=8pt, right=8pt, top=5pt, bottom=5pt, breakable]
\noindent\hyperlink{task:\thepsctr}{\bfseries\color{mainGreen}\fbox{\strut\,$\leftarrow$\,Задача №\csname ztasknum@\thepsctr\endcsname\,}}\quad
\csname zsol@\thepsctr\endcsname
\end{tcolorbox}\par\vspace{0.25cm}%
}%
}
\makeatother
% --- СТРУКТУРНІ ЗАГОЛОВКИ ---
\setcounter{tocdepth}{2}
\newcommand{\chapterTitle}[1]{%
\clearpage\phantomsection\addcontentsline{toc}{section}{#1}%
\sectionTitle{#1}}
\newcommand{\typeTitle}[1]{%
\phantomsection\addcontentsline{toc}{subsection}{\quad #1}%
\par\vspace{0.3cm}
\noindent{\large\bfseries\color{yearOrange}#1}\par\vspace{0.08cm}
\noindent{\color{yearOrange}\rule{\linewidth}{0.4pt}}\par\vspace{0.2cm}}
\newcommand{\subTypeTitle}[1]{%
\phantomsection\addcontentsline{toc}{subsubsection}{\quad\quad #1}%
\par\vspace{0.25cm}
\noindent{\normalsize\bfseries\color{mainGreen}#1}\par\vspace{0.06cm}
\noindent{\color{mainGreen!60}\rule{\linewidth}{0.3pt}}\par\vspace{0.15cm}}
\begin{document}
\chapterTitle{РОЗДІЛ 6. КОМБІНАТОРИКА, ЙМОВІРНІСТЬ, СТАТИСТИКА}
\noindent\small Лише оригінальні завдання НМТ-2026 (перші в кожному блоці), без аналогів. Під кожним~--- кнопка \hyperlink{sol:1}{\bfseries\color{mainGreen}\fbox{\strut\,\footnotesize розв'язок $\rightarrow$\,}}.\par\vspace{0.3cm}
\newpage
\typeTitle{Комбінаторика}
\taskBlock{Комбінаторика за круговою діаграмою --- сесія 23.05, завдання №12}

\noindent
\begin{minipage}[c]{0.55\textwidth}
\zadnum На діаграмі подано інформацію про назви та кількість шкільних гуртків. Макар планує відвідувати два гуртки: один зі спортивних, інший~--- з решти. Скільки різних варіантів вибору має Макар? \nmtyear{2026}
\end{minipage}%
\hfill
\begin{minipage}[c]{0.43\textwidth}
\centering
\begin{tikzpicture}[scale=0.9, thick]
    \def\R{1.7}
    \filldraw[fill=yearOrange!40, draw=white] (0,0) -- (0:\R) arc (0:77.14:\R) -- cycle;
    \filldraw[fill=cyan!30, draw=white] (0,0) -- (77.14:\R) arc (77.14:180:\R) -- cycle;
    \filldraw[fill=mainGreen!30, draw=white] (0,0) -- (180:\R) arc (180:231.43:\R) -- cycle;
    \filldraw[fill=purple!20, draw=white] (0,0) -- (231.43:\R) arc (231.43:257.14:\R) -- cycle;
    \filldraw[fill=yellow!50, draw=white] (0,0) -- (257.14:\R) arc (257.14:360:\R) -- cycle;

    \node at (38.57:1.0) {\small $3$};
    \node at (141.43:1.0) {\small $4$};
    \node at (193:0.9) {\small $1$};
    \node at (244:1.0) {\small $2$};
    \node at (315:1.0) {\small $4$};

    \node at (38.57:2.1) {\scriptsize художні};
    \node at (130:2.2) {\scriptsize танцювальні};
    \node at (200:2.1) {\scriptsize музичні};
    \node[below] at (244:2.0) {\scriptsize STEM};
    \node at (320:2.2) {\scriptsize спортивні};
\end{tikzpicture}
\end{minipage}
\par\vspace{0.2cm}
\answerTable{$36$}{$40$}{$24$}{$14$}{$13$}
\solution{Спортивних гуртків~--- $4$, решти $3+4+1+2=10$. За правилом добутку кількість варіантів $4\cdot 10=40$. Відповідь: \textbf{Б}.}
\vspace{0.3cm}
\taskBlock{Комбінаторика: правило добутку --- сесія 30.05, завдання №7}

\zadtask{У філателістичному магазині є $30$ різних марок з історичної серії та $40$ різних марок із серії <<краєвиди>>. Іван хоче придбати по одній марці з кожної серії. Скільки всього є варіантів такого вибору?}
\answerTable{$70$}{$140$}{$600$}{$1200$}{$1400$}
\solution{За правилом добутку кількість варіантів — добуток кількостей у кожній серії: $30\cdot40=1200$. Відповідь: \textbf{Г} ($1200$).}
\vspace{0.3cm}
\taskBlock{Комбінаторика, кількість способів вибору --- сесія 1.06, завдання №6}

\zadtask{У магазині є $7$ видів подарункових наборів. Скільки всього існує варіантів вибору з них трьох \textit{різних} наборів?}
\answerTable{$21$}{$35$}{$210$}{$343$}{$840$}
\solution{Порядок не важливий, тож це сполучення: $C_7^3=\dfrac{7!}{3!\,4!}=\dfrac{7\cdot6\cdot5}{6}=35$. Відповідь: \textbf{Б}.}
\vspace{0.3cm}
\taskBlock{Комбінаторика: вибір з обов'язковим елементом --- сесія 2.06, завдання №7}

\zadtask{У магазині продають $8$ різних сортів малини, серед яких є сорт <<Альба>>. Покупниця хоче придбати $2$ різні сорти, причому один із них обов'язково має бути <<Альба>>. Скільки в неї є варіантів такого вибору?}
\answerTable{$7$}{$8$}{$14$}{$16$}{$28$}
\solution{Сорт <<Альба>> обов'язково входить, тож залишається вибрати $1$ сорт із решти $8-1=7$. Отже $7$ варіантів. Відповідь: \textbf{А}.}
\vspace{0.3cm}
\taskBlock{Комбінаторика (перестановки) --- сесія 4.06, завдання №11}

\zadtask{Ді-джей складає трек-лист із $6$ різних музичних композицій. Кожен трек має прозвучати лише один раз. Укажіть вираз, за допомогою якого можна визначити загальну кількість можливих послідовностей звучання цих треків.
\vspace{0.2cm}
\par\noindent\textbf{А}\ \ $6 + 5 + 4 + 3 + 2 + 1$
\par\noindent\textbf{Б}\ \ $6$
\par\noindent\textbf{В}\ \ $36$
\par\noindent\textbf{Г}\ \ $6!$
\par\noindent\textbf{Д}\ \ $6^6$
\vspace{0.4cm}}
\solution{Кількість послідовностей із $6$ різних треків (кожен по разу)~--- це число перестановок $P_6=6!$. Відповідь: \textbf{Г}.}
\vspace{0.3cm}
\taskBlock{Комбінаторика: правило добутку (заміна гравця)}

\zadtask{Під час футбольного матчу тренер команди збирається замінити нападником одного із чотирьох захисників. Скільки всього варіантів такої заміни, якщо на лаві запасних два нападники?}
\answerTable{$2$}{$4$}{$6$}{$8$}{$10$}
\solution{Замінити можна будь-кого з $4$ захисників будь-ким із $2$ нападників. За правилом добутку: $4\cdot 2=8$. Відповідь: \textbf{Г}.}
\vspace{0.3cm}
\typeTitle{Теорія ймовірностей}
\taskBlock{Ймовірність за діаграмою опитування --- сесія 1.06, завдання №7}

\noindent
\begin{minipage}[c]{0.54\textwidth}
\zadnum На діаграмі відображено результати анкетування $100$ осіб, які навчаються в університеті, щодо відкриття неподалік нього піцерії. Анкетована особа мала вибрати лише одну з трьох відповідей: так, ні, байдуже. Яка ймовірність того, що навмання обрана особа, яку анкетували, підтримала відкриття піцерії?
\end{minipage}%
\hfill
\begin{minipage}[c]{0.44\textwidth}
\centering
\begin{tikzpicture}[x=1.2cm, y=0.035cm, thick]
    \draw (0,0) -- (0,100);
    \node[rotate=90, font=\small, anchor=south,  yshift=1cm] at (0,50) {Кількість анкетованих осіб};
    \foreach \y in {0,10,...,100} {
        \draw[gray!40, thin] (0,\y) -- (3.5,\y);
        \node[left, font=\small] at (0,\y) {\y};
    }
    \filldraw[fill=cyan!60!blue, draw=none] (0.3,0) rectangle (0.7,80);
    \filldraw[fill=purple!70!black, draw=none] (1.3,0) rectangle (1.7,5);
    \filldraw[fill=mainGreen!80, draw=none] (2.3,0) rectangle (2.7,15);
    \draw (0,0) -- (3.5,0);
    \node[below, font=\small, yshift=-0.1cm] at (0.5,0) {так};
    \node[below, font=\small, yshift=-0.1cm] at (1.5,0) {ні};
    \node[below, font=\small, yshift=-0.1cm] at (2.5,0) {байдуже};
    \node[below, font=\small, yshift=-0.7cm] at (1.5,0) {Варіанти відповідей};
\end{tikzpicture}
\end{minipage}
\par\vspace{0.2cm}
\answerTableTall{$\frac{8}{9}$}{$\frac{1}{5}$}{$\frac{3}{20}$}{$\frac{17}{20}$}{$\frac{4}{5}$}
\solution{Відкриття піцерії підтримали (відповідь «так») $80$ зі $100$ осіб. Шукана ймовірність $P=\dfrac{80}{100}=\dfrac{4}{5}$. Відповідь: \textbf{Д}.}
\vspace{0.3cm}
\taskBlock{Імовірність за списком класу}

\zadtask{У десятому класі за списком $24$ учні. Імовірність того, що вік учня, прізвище якого навмання вибирають зі списку класу, буде меншим за $15$ років, дорівнює $\dfrac{1}{6}$. Скільки всього в класі учнів, вік яких менший за $15$ років?}
\answerTable{$8$}{$4$}{$20$}{$3$}{$6$}
\solution{Кількість учнів, молодших за $15$ років, дорівнює добутку загальної кількості на ймовірність: $24\cdot\dfrac{1}{6}=4$. Відповідь: \textbf{Б}.}
\vspace{0.3cm}
\typeTitle{Статистика (діаграми, графіки)}
\taskBlock{Найбільша різниця за графіками --- сесія 23.05, завдання №14}

\noindent
\begin{minipage}[c]{0.42\textwidth}
\zadnum На рисунку зображено графіки залежності рівня шуму (дБ) від частоти (Гц) для двох джерел. За якої частоти різниця між рівнями шуму джерел є найбільшою? \nmtyear{2026}
\end{minipage}%
\hfill
\begin{minipage}[c]{0.55\textwidth}
\centering
\begin{tikzpicture}
\begin{axis}[
    width=9cm, height=6cm,
    xmode=log, log basis x=10,
    xmin=15, xmax=7000,
    ymin=0, ymax=85,
    ytick={0,20,40,60,80},
    xtick={20,100,200,300,500,1000,2000,5000},
    xticklabels={20,100,200,300,500,1000,2000,5000},
    xlabel={\small частота (Гц)},
    ylabel={\small рівень шуму (дБ)},
    xmajorgrids=true, ymajorgrids=true,
    major grid style={dashed, gray!30},
    axis x line*=bottom, axis y line*=left,
    tick label style={font=\scriptsize},
    log ticks with fixed point,
]
    \addplot[mark=*, thick, mainGreen, mark options={fill=white}] coordinates {
        (20,60) (100,53) (200,62) (300,60) (500,57) (1000,62) (2000,65) (5000,67)
    };
    \addplot[mark=*, thick, yearOrange, mark options={fill=white}] coordinates {
        (20,40) (100,30) (200,20) (300,12) (500,16) (1000,20) (2000,25) (5000,40)
    };
\end{axis}
\end{tikzpicture}
\end{minipage}
\par\vspace{0.2cm}
\answerTable{$1000$}{$5000$}{$200$}{$300$}{$2000$}
\solution{Обчислимо різницю рівнів за кожної частоти: при $300$~Гц вона дорівнює $60-12=48$~дБ, що більше за різниці при інших частотах (напр., $200$: $42$; $1000$: $42$). Найбільша різниця~--- за $300$~Гц. Відповідь: \textbf{Г}.}
\vspace{0.3cm}
\taskBlock{Статистика: кругова діаграма --- сесія 30.05, завдання №4}

\noindent
\begin{minipage}[c]{0.60\textwidth}
\zadnum На круговій діаграмі зображено розподіл туристів за країнами походження, що відвідали певний музей за рік. Визначте відсоток туристів, які прибули з \textit{Канади}.
\end{minipage}%
\hfill
\begin{minipage}[c]{0.38\textwidth}
\centering
\begin{tikzpicture}[scale=1.2, thick]
    \def\R{1.7}
    \filldraw[fill=mainGreen!40, draw=white] (0,0) -- (90:\R) arc (90:-90:\R) -- cycle;
    \node at (0:0.9) {\scriptsize $50\,\%$};
    \node[right] at (0:\R+0.1) {\scriptsize Європа};

    \filldraw[fill=yearOrange!40, draw=white] (0,0) -- (-90:\R) arc (-90:-130.32:\R) -- cycle;
    \node at (-110:1.0) { \scriptsize $11{,}2\,\%$};
    \node[below] at (-110:\R+0.15) {\scriptsize США};

    \filldraw[fill=mainGreen!20, draw=white] (0,0) -- (-130.32:\R) arc (-130.32:-147.6:\R) -- cycle;
    \node[left] at (-139:\R+0.05) {\scriptsize $4{,}8\,\%$ Японія};

    \filldraw[fill=yearOrange!20, draw=white] (0,0) -- (-147.6:\R) arc (-147.6:-160.2:\R) -- cycle;
    \node[left] at (-154:\R+0.05) {\scriptsize $3{,}5\,\%$ Австралія};

    \filldraw[fill=gray!20, draw=white] (0,0) -- (-160.2:\R) arc (-160.2:-270:\R) -- cycle;
    \node at (145:1.0) {\scriptsize $?\,\%$};
    \node[left] at (145:\R+0.05) {\scriptsize Канада};
\end{tikzpicture}
\end{minipage}
\par\vspace{0.2cm}
\answerTable{$25\,\%$}{$28{,}5\,\%$}{$30{,}5\,\%$}{$33\,\%$}{$35\,\%$}
\solution{Сума всіх відсотків дорівнює $100\,\%$. Канада: $100-(50+11{,}2+4{,}8+3{,}5)=100-69{,}5=30{,}5\,\%$. Відповідь: \textbf{В} ($30{,}5\,\%$).}
\vspace{0.3cm}
\taskBlock{Читання стовпчастої діаграми --- сесія 2.06, завдання №6}

\zadtask{На діаграмі відображено розподіл кількості працівників фірми за віком. За даними діаграми визначте різницю між кількістю працівників віком від $20$ до $29$ років та кількістю працівників, яким від $40$ до $49$ років.
\vspace{0.2cm}
\begin{center}
\begin{tikzpicture}[scale=0.9, thick]
    \draw[->, gray!60] (0,0) -- (0,4.8) node[above, font=\scriptsize, text=black] {Кількість};
    \draw[->, gray!60] (0,0) -- (7.4,0) node[right, font=\scriptsize, text=black] {Вік};
    \foreach \v in {2,4,6,8,10,12,14,16,18,20,22,24,26,28,30,32,34,36,38,40,42,44} {
        \pgfmathsetmacro\y{\v/10}
        \draw[gray!20, very thin] (0,\y) -- (6.9,\y);
    }
    \foreach \v in {0,10,20,30,40} {
        \pgfmathsetmacro\y{\v/10}
        \draw[gray!40] (-0.08,\y) -- (6.9,\y);
        \node[left, font=\scriptsize, text=gray!80!black] at (-0.05,\y) {\v};
    }
    \foreach \v in {2,4,6,8,12,14,16,18,22,24,26,28,32,34,36,38,42,44} {
        \pgfmathsetmacro\y{\v/10}
        \draw[gray!40] (-0.06,\y) -- (0.06,\y);
    }
    \foreach \i/\h/\lab in {0/3.6/{20--29}, 1/4.0/{30--39}, 2/2.4/{40--49}, 3/1.6/{50--59}, 4/0.4/{60--70}} {
        \filldraw[fill=cyan!60, draw=white] ({0.6+\i*1.3},0) rectangle ({1.4+\i*1.3},\h);
        \node[below, font=\scriptsize] at ({1.0+\i*1.3},-0.05) {\lab};
    }
\end{tikzpicture}
\end{center}
\par\vspace{0.2cm}}
\answerTable{$6$}{$10$}{$12$}{$16$}{$20$}
\solution{За діаграмою у віці $20$--$29$ років -- $36$ працівників, у віці $40$--$49$ -- $24$. Різниця $36-24=12$. Відповідь: \textbf{В}.}
\vspace{0.3cm}
\taskBlock{Статистика: читання діаграми --- сесія 3.06, завдання №2}

\noindent
\begin{minipage}[c]{0.56\textwidth}
\zadnum На діаграмі відображено інформацію щодо зміни маси (у кг) слона протягом шести місяців із січня до червня цього року в порівнянні з груднем минулого року. Якою стала маса слона в червні, якщо в грудні вона становила $4{,}235$ т?
\end{minipage}%
\hfill
\begin{minipage}[c]{0.42\textwidth}
\centering
\begin{tikzpicture}[x=0.7cm, y=0.035cm, thick]
    \foreach \y in {-60,-40,-20,0,20,40,60,80,100} {
        \draw[gray!60, thin] (0,\y) -- (7,\y);
        \node[left, font=\small] at (0,\y) {\y};
    }
    \draw[gray!60, thin] (0,-65) rectangle (7,100);
    \foreach \x in {0,1,...,6} {
        \draw[gray!60, thin] (\x,-65) -- (\x,100);
    }
    \draw[thick] (0,0) -- (7,0);
    \node[rotate=90, font=\small, yshift=0.8cm] at (0,20) {Зміна маси слона, кг};
    \foreach \x/\mon in {0/грудень, 1/січень, 2/лютий, 3/березень, 4/квітень, 5/травень, 6/червень} {
        \node[below, rotate=90, anchor=east, font=\small, xshift=-0.1cm] at (\x,-65) {\mon};
    }
    \node[below, font=\small, yshift=-2.2cm] at (3.5,-65) {Місяць};
    \foreach \x/\y in {0/0, 1/-60, 2/-40, 3/0, 4/20, 5/80, 6/60} {
        \filldraw[black] (\x,\y) circle (2pt);
    }
\end{tikzpicture}
\end{minipage}
\par\vspace{0.2cm}
\answerTable{$4305$ кг}{$4295$ кг}{$4235$ кг}{$4275$ кг}{$4215$ кг}
\solution{Маса в грудні $4{,}235$ т $= 4235$ кг. За діаграмою зміна маси у червні порівняно з груднем дорівнює $+60$ кг. Отже маса в червні: $4235 + 60 = 4295$ кг. Відповідь: \textbf{Б}.}
\vspace{0.3cm}
\taskBlock{Читання діаграми (ціни) --- сесія 4.06, завдання №2}

\noindent
\begin{minipage}[t]{0.52\textwidth}
\vspace{0pt}
\zadnum На діаграмі відображено інформацію про ціну пакета цукерок «Чудасія», «Каштанчики», «Шокос», «Смакс» і «Няммі» в деякому магазині. Якщо в Катрусі є лише $200$~грн, то із зазначених цукерок вона може придбати
\vspace{0.3cm}
\par\noindent\textbf{А}\ \ лише «Каштанчики» або «Шокос»
\par\noindent\textbf{Б}\ \ лише «Чудасія»
\par\noindent\textbf{В}\ \ лише «Чудасія», «Смакс» або «Няммі»
\par\noindent\textbf{Г}\ \ будь-який один пакет цукерок
\par\noindent\textbf{Д}\ \ лише «Каштанчики», «Смакс» або «Няммі»
\end{minipage}%
\hfill
\begin{minipage}[t]{0.48\textwidth}
\vspace{0pt}
\centering
\begin{tikzpicture}[scale=0.85, thick]
    \draw[gray!60] (0,0) -- (6.5,0);
    \draw[gray!60] (0,0) -- (0,4.0);
    \node[rotate=90, font=\small] at (-1.2, 2.0) {Ціна (грн за пакет)};
    \foreach \v/\y in {0/0, 50/0.5, 100/1.0, 150/1.5, 200/2.0, 250/2.5, 300/3.0, 350/3.5} {
        \draw[gray!20, thin] (0,\y) -- (6.5,\y);
        \node[left, font=\scriptsize, text=gray!80!black] at (0,\y) {\v};
    }
    \filldraw[fill=yearOrange!70, draw=white] (0.4,0) rectangle (1.2, 1.75);
    \filldraw[fill=mainGreen!80, draw=white] (1.6,0) rectangle (2.4, 2.5);
    \filldraw[fill=brown!70, draw=white] (2.8,0) rectangle (3.6, 3.25);
    \filldraw[fill=cyan!70, draw=white] (4.0,0) rectangle (4.8, 1.9);
    \filldraw[fill=purple!60, draw=white] (5.2,0) rectangle (6.0, 2.0);
    \node[below, rotate=50, anchor=north east, font=\scriptsize, inner sep=2pt] at (1.0, -0.1) {«Чудасія»};
    \node[below, rotate=50, anchor=north east, font=\scriptsize, inner sep=2pt] at (2.2, -0.1) {«Каштанчики»};
    \node[below, rotate=50, anchor=north east, font=\scriptsize, inner sep=2pt] at (3.4, -0.1) {«Шокос»};
    \node[below, rotate=50, anchor=north east, font=\scriptsize, inner sep=2pt] at (4.6, -0.1) {«Смакс»};
    \node[below, rotate=50, anchor=north east, font=\scriptsize, inner sep=2pt] at (5.8, -0.1) {«Няммі»};
\end{tikzpicture}
\end{minipage}
\solution{За діаграмою ціни: «Чудасія» $175$, «Каштанчики» $250$, «Шокос» $325$, «Смакс» $190$, «Няммі» $200$~грн. На $200$~грн вистачає лише на ті, що коштують не більше $200$: «Чудасія», «Смакс» або «Няммі». Відповідь: \textbf{В}.}
\vspace{0.3cm}
\taskBlock{Читання діаграми (результати опитування) --- сесія 5.06, завдання №2}

\noindent
\begin{minipage}[c]{0.54\textwidth}
\zadnum На рисунку відображено результати опитування пасажирів щодо якості перевезень (у балах за шкалою $1$--$10$). Визначте кількість опитаних, які поставили оцінку, \textit{вищу за} $8$ балів.
\end{minipage}%
\hfill
\begin{minipage}[c]{0.44\textwidth}
\centering
\begin{tikzpicture}[scale=0.42, thick]
    \draw[->, gray!60] (0,0) -- (0,9) node[above, font=\scriptsize] {};
    \draw[->, gray!60] (0,0) -- (7.5,0) node[right, font=\scriptsize, text=black] {бали};
    \foreach \v/\yy in {25/1,50/2,75/3,100/4,125/5,150/6,175/7,200/8,225/9} {
        \draw[gray!20] (0,\yy) -- (7,\yy);
    }
    \foreach \v/\yy in {50/2,100/4,150/6,200/8} {\node[left,font=\tiny, text=gray!80!black] at (0,\yy){\v};}

    \draw[mainGreen, thick] (1,3)--(2,4)--(3,4)--(4,8)--(5,7)--(6,3)--(7,1);
    \foreach \xx/\yy in {1/3,2/4,3/4,4/8,5/7,6/3,7/1} {
        \filldraw[fill=white, draw=mainGreen, thick] (\xx,\yy) circle (3.5pt);
    }
    \foreach \xx/\lab in {1/4,2/5,3/6,4/7,5/8,6/9,7/10} {\node[below,font=\tiny] at (\xx,0){\lab};}
\end{tikzpicture}
\end{minipage}
\par\vspace{0.2cm}
\answerTable{$19$}{$100$}{$200$}{$275$}{$475$}
\solution{Оцінка вища за $8$ балів~--- це $9$ і $10$ балів. За графіком їх поставили $75$ і $25$ опитаних. Разом $75+25=100$. Відповідь: \textbf{Б}.}
\vspace{0.3cm}
\printAllSolutions
\chapterTitle{ВІДПОВІДІ}
\noindent\small Відповіді наведено у тому ж порядку, що й завдання.\par\vspace{0.3cm}
\ansType{Завдання з вибором однієї відповіді (А--Д)}
\begin{multicols}{3}\noindent
\ansitem{Б}
\ansitem{Г}
\ansitem{Б}
\ansitem{А}
\ansitem{Г}
\ansitem{Г}
\ansitem{Д}
\ansitem{Б}
\ansitem{Г}
\ansitem{В}
\ansitem{В}
\ansitem{Б}
\ansitem{В}
\ansitem{Б}
\end{multicols}
\end{document}
